% HistoWeave — Bioinformatics Original Paper
% Template: Oxford University Press / Bioinformatics
\documentclass[]{bioinfo}

\usepackage[utf8]{inputenc}
\usepackage{amsmath,amssymb}
\usepackage{booktabs}
\usepackage{graphicx}
\usepackage{xcolor}
\usepackage{hyperref}
\usepackage{multirow}
\usepackage{siunitx}

% -- Metadata --
\copyrightyear{2026}
\pubyear{2026}
\access{Advance Access Publication Date: Day Month Year}
\appnotes{Original Paper}

\begin{document}
\firstpage{1}

\subtitle{Gene Expression}

\title[HistoWeave]{HistoWeave: evidence-governed method decisions with structured
abstention for spatial transcriptomics}

\author[Author~A \textit{et~al}.]{
Author~A\,$^{1,*}$,
Author~B\,$^{1}$,
Author~C\,$^{2}$
and Author~D\,$^{1}$
}

\address{
$^{1}$Department of XXX, University of XXX, City, Country \\
$^{2}$Department of XXX, University of XXX, City, Country
}

\corresp{$^\ast$To whom correspondence should be addressed.}

\history{Received on XXXXX; revised on XXXXX; accepted on XXXXX}

\editor{Associate Editor: XXXXX}

\abstract{
\textbf{Motivation:}
Spatial transcriptomics methods yield substantially different domain
partitions depending on data characteristics, spatial weight
parameterisation, and---critically---whether the true domain count is
provided as an oracle input. Despite growing method arsenals, no
systematic framework exists that (i)~rejects cross-task or circularly
validated evidence before aggregation, (ii)~returns set-valued
method decisions rather than a forced single winner, and (iii)~abstains
from personalisation when held-out evidence is absent or negative.

\textbf{Results:}
We present HistoWeave, an open-source evidence-governed decision
protocol for spatial transcriptomics method selection. HistoWeave
enforces executable admissibility gates on task semantics, ground-truth
kind, and oracle-\(K\) status, then emits one of four auditable actions:
\texttt{personalised\_set}, \texttt{global\_default},
\texttt{evidence\_required}, or \texttt{abstain}. The oracle-\(K\)
DLPFC track contains 20 method configurations; a separate oracle-\(K\)
external panel contains 15 prespecified methods (13 with finite results)
on five datasets spanning Visium, Visium~HD, Xenium, and MERFISH.
Removing oracle-\(K\) protection reduces SpaGCN ARI by up to 55\% on an
individual DLPFC slice. External five-dataset LOOCV ties, but does not
beat, the global-best comparator (both mean regret 0.006). A
prospectively specified in-repository test on six unseen breast-cancer
sections also fails its margin (mean regret 0.131, 95\% CI
\([0.034, 0.236]\); threshold 0.02); because it lacks an independent
public timestamp and uses oracle \(K\), we treat it as a negative stress
test rather than confirmatory validation of non-oracle personalisation.
Finally, an 11-case adversarial audit yields zero incompatible-evidence
admissions. HistoWeave preserves these claim boundaries and negative
results as machine-readable artefacts.

\textbf{Availability:}
HistoWeave is available at
\url{https://github.com/HERRY423/Histoweave}
under the BSD-3-Clause licence (PyPI: \texttt{histoweave-spatial}).
Frozen figures, benchmark data, and reproduction scripts are in
\texttt{submission\_freeze\_v2/}.
DOI: \url{https://doi.org/10.5281/zenodo.21586217}

\textbf{Contact:} \href{mailto:xxx@xxx.edu}{xxx@xxx.edu}

\textbf{Supplementary information:} Supplementary data are available
at \textit{Bioinformatics} online.
}

\maketitle

% =====================================================================
\section{Introduction}

Spatial transcriptomics technologies---including Visium, Visium~HD,
MERFISH, Slide-seqV2, and Xenium---now routinely measure gene expression
while preserving tissue architecture
\citep{maynard2021,stahl2016,chen2022}. The explosion of analysis
methods (STAGATE, BayesSpace, SpaGCN, GraphST, BANKSY, and dozens of
general-purpose alternatives) creates a practical crisis for biologists:
\textit{which method should I apply to my data, and how much should I
trust the result?}

Current benchmark studies typically compare methods on a single
platform, provide the true domain count (\(K\)) as an oracle input, and
report a single ``best'' method via averaged metrics
\citep{li2023benchmarking}. This practice conceals three systematic
biases. First, \textbf{oracle-\(K\) inflation}: methods that require
\(K\) as input benefit from leaked ground-truth information, inflating
reported accuracy by an uncontrolled margin. Second, \textbf{ground-truth
semantic mismatch}: spatial priors may help when ground-truth labels
represent contiguous tissue layers but harm performance when labels
denote transcriptomic cell types scattered across the tissue. Third,
\textbf{cross-task contamination}: aggregating evidence across
incompatible analysis tasks (e.g., domain detection vs.\ cell-type
annotation) produces rankings that cannot be acted upon.

We introduce HistoWeave, an evidence-governed decision protocol that
addresses these shortcomings. Rather than selecting a universal winner,
HistoWeave produces an auditable \texttt{DecisionCard} that records
which evidence was admitted, which was rejected (and why), and whether
the data justify a personalised method set, a global default, or
abstention. The protocol is implemented as open-source software with
typed plugins for 80+ analysis methods, containerised SOTA backends,
multi-objective Pareto analysis, and a federated evidence network.

Our contributions are:
\begin{enumerate}
  \item An \textbf{executable evidence-admissibility protocol} that
        machine-checks task semantics, ground-truth kind, and oracle-\(K\)
        status before any score influences a decision (Section~\ref{sec:protocol}).
  \item Quantification of \textbf{oracle-\(K\) leakage}: SpaGCN loses up
        to 55\% relative ARI on DLPFC when deprived of the true domain
        count (Section~\ref{sec:oracle-k}).
  \item A \textbf{coverage-stratified cross-platform evaluation} that
        reports the DLPFC and external method panels separately and finds
        no personalisation advantage over the global comparator
        (Section~\ref{sec:crossplatform}).
  \item A \textbf{fail-closed validation chain}: an 11-case adversarial
        audit rejects every incompatible evidence object, while a negative
        one-shot external stress test remains a first-class artefact
        (Section~\ref{sec:independent-test}).
\end{enumerate}

% =====================================================================
\section{Methods}

\subsection{Design principles}

HistoWeave is guided by four principles: (i)~every published ARI must
carry a protocol tag declaring \(K\) policy, seed policy, and
ground-truth kind; (ii)~negative results are first-class---personalisation
that fails the global gate must stay \texttt{global\_default};
(iii)~cross-task, circular, and proxy-domain evidence is rejected rather
than softly down-weighted; and (iv)~reference artefacts are frozen paths
whose re-runs must match or bump the protocol version.

\subsection{Evidence-governed decision protocol}
\label{sec:protocol}

For a declared task \(\tau\) and query dataset \(q\), the protocol
applies an ordered rule chain:

\begin{enumerate}
  \item \textbf{Admissibility gate.} Reject evidence whose task, ground-truth
        kind, or oracle-\(K\) status is incompatible with \(\tau\).
        Cross-modal partitions (RNA vs.\ protein vs.\ chromatin) are
        isolated; cluster-proxy labels are excluded.
  \item \textbf{Candidate generation.} Extract target-free features from
        \(q\), retrieve the \(k\) nearest reference datasets from the
        landscape knowledge base, and rank method configurations by
        similarity-weighted ARI. The recommender's confidence field is an
        uncalibrated rank-support heuristic, not a probability.
  \item \textbf{Global-default comparison.} If the local proxy does not
        beat the global-default proxy, return \texttt{global\_default}.
  \item \textbf{Held-out validation gate.} If independent grouped held-out
        validation is missing, return \texttt{evidence\_required} even
        when the local proxy appears favourable.
  \item \textbf{Pareto intersection.} When matched multi-objective data are
        supplied, return only non-dominated configurations; different
        spatial-weight parameterisations (\texttt{method@sw0.0} vs.\
        \texttt{method@sw0.8}) are distinct decisions.
  \item \textbf{Decision.} Return \texttt{personalised\_set} if all gates
        pass; \texttt{abstain} when no task-valid evidence remains.
\end{enumerate}

The output is a versioned \texttt{DecisionCard}---a machine-readable
record of evidence roles, contract checks, claim boundaries, and the
selected action. Default thresholds are serialised with the card;
changing them creates a different decision policy.

\subsection{Information-theoretic Spatial Utility Score (ISUS)}

We define ISUS as the ratio of conditional to marginal mutual
information:
\begin{equation}
\text{ISUS} = \frac{I(D; S \mid E)}{I(D; E)}
\label{eq:isus}
\end{equation}
where \(D\) are domain labels, \(S\) spatial coordinates, and \(E\)
expression features. The numerator captures domain information that
spatial coordinates provide \emph{beyond} expression; the denominator
normalises by what expression alone provides. Both terms use the
Ross~(2014) \(k\)NN estimator. ISUS requires trusted labels and is
therefore a post-hoc diagnostic, not a pre-execution selector.

\subsection{Multi-objective Pareto analysis}

Method configurations are evaluated on up to four objectives: accuracy
(ARI, maximise), speed (seconds, minimise), memory (GB, minimise), and
robustness (bootstrap ARI CI width, minimise). The Pareto layer reports
the non-dominated frontier---the set of configurations that no other
configuration beats on \emph{every} axis. Among frontier members,
the \emph{knee} is the point closest to the ideal corner in normalised
objective space.

\subsection{Software architecture}

HistoWeave is implemented in Python~3.11+ as a typed plugin framework.
The \texttt{SpatialTable} data model wraps \texttt{SpatialData} and
\texttt{AnnData}, enforcing provenance tracking via structured
\texttt{uns["provenance"]} records. Method plugins declare maturity
tiers (\textit{experimental} $\to$ \textit{beta} $\to$
\textit{production} $\to$ \textit{contract\_validated} $\to$
\textit{validated}). SOTA backends (STAGATE, GraphST, BayesSpace) use
fail-closed container isolation: if the external environment is
unavailable, the system raises an error rather than substituting a
trivial fallback.

% =====================================================================
\section{Results}

\subsection{DLPFC benchmark: 20 configurations under one protocol}
\label{sec:dlpfc}

We benchmarked 20 method configurations---five scikit-learn algorithms
each evaluated at three spatial weights
(\(w \in \{0.0, 0.3, 0.8\}\)), plus five spatial-aware methods
(STAGATE, BayesSpace, SpaGCN, GraphST, BANKSY)---on five DLPFC slices
from \citet{maynard2021}, each run with three random seeds (42, 1, 2)
and truth-derived \(K\) (oracle track).

STAGATE achieved the highest grand mean ARI
(0.353), followed by BayesSpace (0.326) and SpaGCN (0.317)
(Table~\ref{tab:dlpfc}). Spatial-aware methods outperformed sklearn
baselines on average (mean ARI 0.287 vs.\ 0.186), but the best sklearn
configuration (\texttt{gaussian\_mixture@sw0.8}, ARI 0.254) remained
competitive on several slices.

No single method universally dominated. Slice~151669 (\(K=8\)) proved
uniformly difficult---no method exceeded ARI 0.25---while
slice~151670 (\(K=5\)) showed atypical rankings, with
\texttt{spectral@sw0.3} (ARI 0.346) outperforming all SOTA methods.
BayesSpace exhibited the lowest seed variance (std $\leq$ 0.044),
while STAGATE showed the highest on certain slices (std up to 0.076).

\begin{table}[!t]
\procemark
\caption{Grand mean ARI over 5 DLPFC slices $\times$ 3 seeds (oracle-$K$
track). Top 10 of 20 configurations shown; full matrix in Supplementary
Table~S1.\label{tab:dlpfc}}
\begin{tabular}{@{}llr@{}}
\toprule
Method & Family & Mean ARI \\
\midrule
STAGATE         & spatial-aware & 0.353 \\
BayesSpace      & spatial-aware & 0.326 \\
SpaGCN          & spatial-aware & 0.317 \\
GMM\texttt{@sw0.8} & sklearn   & 0.254 \\
spectral\texttt{@sw0.8} & sklearn & 0.243 \\
kmeans\texttt{@sw0.3}  & sklearn  & 0.235 \\
agglom.\texttt{@sw0.8} & sklearn  & 0.230 \\
spectral\texttt{@sw0.3} & sklearn & 0.227 \\
birch\texttt{@sw0.8}   & sklearn  & 0.226 \\
BANKSY          & spatial-aware & 0.223 \\
\bottomrule
\end{tabular}
\end{table}

\subsection{Oracle-$K$ leakage quantification}
\label{sec:oracle-k}

Published spatial transcriptomics benchmarks routinely provide the true
domain count \(K\) as input, conflating method performance with leaked
ground-truth information. We ran SpaGCN and STAGATE in dual-track mode:
oracle-\(K\) (truth-derived) versus estimate-\(K\) (silhouette-based).

SpaGCN's oracle-\(K\) mean ARI (0.299) dropped to 0.237 under silhouette
estimation ($\Delta = -0.062$; Table~\ref{tab:oracle-k}). On the most
affected slice (151673), SpaGCN collapsed from ARI 0.418 (oracle
\(K=7\)) to 0.186 (estimated \(K=2\)), a 55\% relative loss.
STAGATE showed a smaller but consistent drop ($\Delta = -0.013$). All
blind estimators (silhouette, spatial silhouette, ensemble) converged
toward \(K=2\), failing to recover the true 5--8 cortical layers
(Table~\ref{tab:oracle-k}).

\begin{table}[!t]
\procemark
\caption{Oracle-$K$ leakage: mean ARI under oracle vs.\ estimated $K$
(seed 42, DLPFC).\label{tab:oracle-k}}
\begin{tabular}{@{}lrrr@{}}
\toprule
Method & Oracle-$K$ ARI & Estimate ARI & $\Delta$ \\
\midrule
SpaGCN  & 0.299 & 0.237 & $-$0.062 \\
STAGATE & 0.232 & 0.219 & $-$0.013 \\
\midrule
\multicolumn{4}{@{}l}{\textit{Worst case: SpaGCN slice 151673}} \\
Oracle ($K$=7)  & \multicolumn{2}{c}{0.418} & \\
Silhouette ($K$=2) & \multicolumn{2}{c}{0.186} & $-$0.232 \\
\bottomrule
\end{tabular}
\end{table}

This finding demonstrates that dual-track reporting (oracle vs.\
estimated \(K\), with ARI under both) is the minimal audit trail for
any SOTA claim that includes an oracle-\(K\) number.

\subsection{Cross-platform external validation}
\label{sec:crossplatform}

We assembled five external datasets spanning four platforms and multiple
tissues (Table~\ref{tab:external}). All datasets carry spatial-region
ground truth (anatomical, pathology, or manual spatial-region labels;
never cell-type predictions). The panel prespecified 15 methods, but
\texttt{spatialde\_kmeans} and \texttt{nnsvg\_kmeans} produced no finite
ARI; the common finite panel therefore contains 13 methods. Except for
BANKSY, the five spatial-aware SOTA methods in the DLPFC table were not
run across this external panel. All external scores in this section use
truth-derived \(K\) and are consequently oracle-track evidence.

\begin{table}[!t]
\procemark
\caption{External validation datasets and best-performing
methods.\label{tab:external}}
\begin{tabular}{@{}llllr@{}}
\toprule
Dataset & Platform & Tissue & Best method & ARI \\
\midrule
CRC    & Visium HD     & Human colon  & spectral & 0.676 \\
Lung   & Xenium        & Human lung   & mean\_shift & 0.183 \\
Ovarian & Xenium Prime & Human ovary  & spectral & 0.463 \\
Mouse brain & Visium v2 & Mouse brain & spectral & 0.627 \\
MERFISH & MERFISH      & Mouse brain  & spectral & 0.282 \\
\bottomrule
\end{tabular}
\end{table}

Using leave-one-out cross-validation (LOOCV) across the five held-out
datasets, the $k$NN recommender achieved top-1 accuracy of 0.80 and mean
selection regret of 0.006 ARI---tying, but not beating, the global-best
baseline (spectral clustering). This result reduces regret by 97.5\%
relative to random selection but does not justify personalised deployment
(Table~\ref{tab:loocv}).

\begin{table}[!t]
\procemark
\caption{Recommendation LOOCV summary (5 external
datasets).\label{tab:loocv}}
\begin{tabular}{@{}lr@{}}
\toprule
Metric & Value \\
\midrule
$n$ queries & 5 \\
Top-1 accuracy & 0.80 \\
Mean selection regret & 0.006 \\
Global-best regret & 0.006 \\
Random-choice regret & 0.234 \\
Regret reduction vs.\ random & 97.5\% \\
Regret reduction vs.\ global & 0.0\% \\
\bottomrule
\end{tabular}
\end{table}

Critically, the broader study-grouped evaluation (20~independent
queries) showed that always-personalised regret (0.047) \emph{exceeded}
always-global regret (0.029). HistoWeave's selective protocol correctly
detected this and chose zero personalisation coverage, operationalised as
\texttt{global\_default} at every operating point. This is abstention
from personalisation, not the distinct \texttt{abstain} action used when
no task-valid evidence remains.

\subsection{Prospectively specified external stress test (negative result)}
\label{sec:independent-test}

To test whether the frozen spectral policy generalises beyond
brain-tissue datasets, we recorded an in-repository protocol before the
local outcome run for six breast-cancer Visium sections from
\citet{wu2021} (Zenodo DOI 10.5281/zenodo.4739739). The study was absent
from the audited training and model-selection sources. The protocol
locked the policy (\texttt{spectral}), seven-method comparator panel,
oracle-\(K\) task contract, and a success margin of 0.02~ARI regret.
The repository record has no independent public timestamp; accordingly,
we do not describe this analysis as publicly preregistered or use it to
unlock the non-oracle personalisation action.

The frozen spectral policy achieved mean regret 0.131 (95\% bootstrap
CI: $[0.034, 0.236]$), exceeding the 0.02 threshold by an order of
magnitude. Spectral was the top-1 method on only 2 of 6 sections (33\%).
The analysis returned \texttt{independent\_test\_fail} and
retained this negative result as a scientific artefact rather than
discarding it.

This result demonstrates two things: (i)~a policy optimised on brain
cortex does not transfer to breast tumour microenvironments, and
(ii)~the negative endpoint is consistent with retaining the global
fallback. Because this test uses oracle \(K\), it is descriptive stress
evidence and not admissible held-out evidence for non-oracle personalised
deployment. The test cohort remains excluded from training.

\subsection{Scalability}

We profiled 30 HistoWeave pipeline methods on synthetic datasets
from \num{1000} to \num{1000000} cells ($\times$ 2000 genes), recording
wall-clock time and peak RSS within a 64~GB memory cap and 1800~s
timeout. 10 of 30 methods reached the 1M-cell scale; 20 were
terminated by memory limits at intermediate scales. The fastest method
at 1M cells completed in 40.7~s; time complexity exponents ranged from
sub-linear (QC operations) to super-quadratic (spatial graph methods).

% =====================================================================
\section{Discussion}

HistoWeave contributes an \emph{evidence-governed decision protocol}
rather than a new clustering algorithm. Its core proposition is that
method selection in spatial transcriptomics should be treated as a
structured decision problem with explicit admissibility checks, fallback
rules, and abstention---not as a ranking exercise that always produces
a winner.

\textbf{Oracle-$K$ as a hidden confounder.}
Our dual-track analysis reveals that the common practice of providing
the true domain count inflates reported ARI by up to 55\% for individual
method--slice combinations. This effect is not a minor calibration
issue: on DLPFC slice 151673, SpaGCN's estimated-\(K\) ARI (0.186) would
rank it below several non-spatial baselines. We argue that dual-track
\(K\) reporting should become a standard requirement for spatial
transcriptomics benchmarks.

\textbf{Ground-truth semantics matter.}
All external outcomes reported here use spatial-region ground truth, so
they support comparisons within that task only. They do not test, and
cannot support a claim about, performance on transcriptomic cell-type
labels. HistoWeave nevertheless records ground-truth semantics because
mixing these tasks during evidence aggregation would make a method
decision uninterpretable.

\textbf{Negative results as evidence.}
The prospectively specified in-repository breast-cancer stress test
failed its margin by an order of magnitude. We retain it as a first-class
scientific artefact while explicitly distinguishing it from a publicly
timestamped preregistration and from admissible non-oracle deployment
evidence. The test cohort remains excluded from future training.

\textbf{Comparison with existing tools.}
HistoWeave complements rather than replaces SpatialData (data model and
I/O) and Squidpy (spatial analysis grammar). Its irreplaceable layer is
the evidence-governed decision workflow: task/ground-truth admissibility
gates, structured abstention, fail-closed SOTA backends, and study-level
generalisation statistics. Open Problems in Single-Cell Analysis provides
general benchmarking infrastructure but does not enforce task-specific
evidence admissibility or produce set-valued decisions with abstention.

\textbf{Limitations.}
The current external landscape contains only 5--10 independent
validation units, limiting the statistical power of recommendation
LOOCV. Absolute ARI values on non-DLPFC tissues remain moderate
(0.18--0.68), partly reflecting genuine difficulty and partly the
limitation of the external method panel. External evaluation is
oracle-\(K\), and four of five DLPFC SOTA families are missing from the
five-dataset panel; no cross-panel SOTA superiority claim is made. The
ISUS gain calibration on five DLPFC slices is underpowered (Spearman
$\rho \approx -0.30$), and ISUS remains a post-hoc descriptor. Future
work will expand real independent study units beyond 15, incorporate
non-oracle \(K\) estimators as standard infrastructure, complete an
aligned SOTA external panel, and obtain a public timestamp before the
next untouched prospective validation.

% =====================================================================
\section{Availability and Implementation}

HistoWeave is available at
\url{https://github.com/HERRY423/Histoweave} under the BSD-3-Clause
licence. The Python package \texttt{histoweave-spatial} is distributed
via PyPI. The P0 submission freeze (manuscript, evidence-admission audit,
method-coverage ledger, source-artifact hashes, and reproduction script)
is in \texttt{submission\_freeze\_v2/}. Frozen benchmark
data, prepared dataset scripts, and environment locks for SOTA methods
are tracked in the repository. Raw spatial transcriptomics datasets are
not redistributed; preparation scripts document how each public dataset
is transformed. The Wu~et~al.\ 2021 test cohort remains excluded from
training.
Archive DOI: \url{https://doi.org/10.5281/zenodo.21586217}

\section*{Acknowledgements}

% Disclose AI assistance if retained in the submitted work.

\section*{Funding}

This work was supported by XXX [grant number XXX].

\bibliographystyle{natbib}
% \bibliography{references}

\begin{thebibliography}{}

\bibitem[Chen \textit{et~al.}(2022)]{chen2022}
Chen,~A. \textit{et~al.} (2022)
Spatiotemporal transcriptomic atlas of mouse organogenesis using DNA
nanoball-patterned arrays.
\textit{Cell}, \textbf{185}(10), 1777--1792.

\bibitem[Li \textit{et~al.}(2023)]{li2023benchmarking}
Li,~B. \textit{et~al.} (2023)
Benchmarking spatial and single-cell transcriptomics integration methods
for transcript distribution prediction and cell type deconvolution.
\textit{Nature Methods}, \textbf{19}(6), 662--670.

\bibitem[Maynard \textit{et~al.}(2021)]{maynard2021}
Maynard,~K.R. \textit{et~al.} (2021)
Transcriptome-scale spatial gene expression in the human dorsolateral
prefrontal cortex.
\textit{Nature Neuroscience}, \textbf{24}(3), 425--436.

\bibitem[Palla \textit{et~al.}(2022)]{palla2022squidpy}
Palla,~G. \textit{et~al.} (2022)
Squidpy: a scalable framework for spatial omics analysis.
\textit{Nature Methods}, \textbf{19}(2), 171--178.

\bibitem[Ross(2014)]{ross2014}
Ross,~B.C. (2014)
Mutual information between discrete and continuous data sets.
\textit{PLoS ONE}, \textbf{9}(2), e87357.

\bibitem[Stahl \textit{et~al.}(2016)]{stahl2016}
St{\aa}hl,~P.L. \textit{et~al.} (2016)
Visualization and analysis of gene expression in tissue sections by
spatial transcriptomics.
\textit{Science}, \textbf{353}(6294), 78--82.

\bibitem[Wu \textit{et~al.}(2021)]{wu2021}
Wu,~S.Z. \textit{et~al.} (2021)
A single-cell and spatially resolved atlas of human breast cancers.
\textit{Nature Genetics}, \textbf{53}(9), 1334--1347.

\end{thebibliography}

\end{document}
